% Online Appendix to Appendix D (app:corrections). Item-level detail of the corrections listed in
% Table tab:corrections_index: defect, symptom in the shipped release, correction, and the
% pre-rerun effect measured by the audit on the shipped inputs (not a result of the paper).
% Sources: the audit's fix logs and reviews, git log of branch fix/allied-base-paper, README
% changelog, and the corrected code in R/, scripts/, config/.
\section{Corrections to the shipped release: item detail}\label{oa:corrections}

This appendix gives, for each item of Table~\ref{tab:corrections_index}, the defect, its symptom in
the shipped release of September 2026, the correction and, where the audit measured it, the
\emph{pre-rerun effect}: the change the correction makes when applied alone to the theme tables and
inputs saved with the shipped release. Pre-rerun effects isolate one correction on inputs that
still embed the other defects; they are not results of the paper.

\paragraph{Aggregation (A1--A7).} A1: an unobserved category kept its weight in the denominator and
counted as zero, and a cell with no observed category scored zero; 1{,}809 of 2{,}773 headline ES
Reserves rows were unobserved. The available-case rule alone changes 1{,}682 Upstream cells and
raises mean Upstream ES by 0.139 (Oil), 0.112 (Gas) and 0.104 (Coal). Because the same commit gives
Reserves a zero rule (the reserve tables list every material holder, so absence means none), the two
changes offset on the shipped tables; the rule matters for categories without a missing-data rule
and for the PSI (B4). A2: a zero or global-average rule filled groups in which no country was
observed; it now fills only groups with an observed country and reports the rest. A3:
\texttt{any(isTRUE(x))} is \texttt{FALSE} for vectors longer than one, so the diagnostics reported
11.9 and 29.5 percent of ES and EO category scores as imputed against 37.0 and 69.2 percent in fact;
the final run has 38.4 and 45.0 percent. A4: Energy Prices was configured on a variable no theme
produced and dropped from both pillars without warning, because the guard compared the join keys with
themselves; it re-enters ten Upstream cells as a world price series that shifts cell levels and
leaves within-cell order unchanged (Spearman 1.000 in every varying cell). A second guard now warns
when a theme's own composite is scored in place of a configured overall. A5: an EO Reserves weight
of 2 for a category EO does not have let resource-potential rows create 88 fully imputed EO cells
(the Holy See, the Pitcairn Islands); EO falls from 6{,}923 to 6{,}835 headline cells. A6: the
variable-contribution table divided by the number of components including missing ones; both
decompositions now reconcile to the index (largest error $1.1\times10^{-16}$). A7: a development
commit made \texttt{bind\_rows()} read a named list of tables as one row; a rerun from that revision
would have removed every overall-based category.

\paragraph{Country universe (C1--C3).} C1: the standardiser attached ISO3 codes but did not rename
rows, so the shipped ES, EO, Strategic and Policy tables carried 272 strings for 237 economies, with
30 codes under two or three spellings (Viet Nam and Vietnam, Korea and South Korea, the two Congos).
Each spelling took its own rank in every S-curve with a partly observed, often imputed profile.
Every variant now maps to the WDI name of its code \citep{worldbank2026wdi}; territories with their
own code stay separate. C2: the EI ``USSR'' series, a CIS aggregate, was coded RUS, ``Northern
Cyprus'' CYP, and a pseudo-country ``Global'' carried readiness values (128 rows); aggregates,
residuals and historical entities are now dropped, not merged into successors. C3: the 2025
critical-minerals Comtrade pulls hold three reporters against about 130 in 2024; a vintage needs 50.

\paragraph{Energy Security themes (E1--E12).} E1: graphite and rare earths fell out of every join
because the IEA and EI name them differently; with harmonised names the rare earths carry about 79
percent of Wind's Mineral Supply weight. E2: the stage recode mapped refining to the Upstream under
both branches; refining now scores the Midstream. E3: listed non-producers (Germany, Japan) tied at
0.441 just below the Congo at 0.473 in Batteries Upstream while unlisted economies scored zero; zero
supply now maps to zero for every economy and suppliers rank above it. E4: the reserves roll-up
averaged only over listed minerals, so Chile scored 1.000 on Wind Reserves on copper alone. E5: BNEF
onshore and offshore wind were never recoded, so Wind had no Consumption category; growth was ranked
across technologies; Viet Nam's rows fall from 160 to 63 once each country receives one regional
value. E6: Minerals Trade kept only the market-share index; it is now the mean of a demand-weighted
market share and a concentration index $\Scurve(-\HHI_m)$. E7: critical-minerals investment was
copied in full to Batteries and EV Upstream (79.8 each; now 9.2 and 70.6), and stage completion
manufactured zero Investment cells in Nuclear Upstream and Midstream and Geothermal Upstream and
Midstream, which made those cells constant across countries and gave Nuclear Upstream the maximum
security need for every country. E8: electrolyser capacity was filed outside the headline
technologies. E9: mineral Upstream production depth was an unweighted mean of different minerals'
kilotonnes; it is now a demand-weighted share of world mined supply, with growth as an absolute
change (EO moves by up to 0.067 in six cells, within-cell Spearman 0.955--0.992). E10: a placeholder
region dropped the eight non-EU European EI economies from EV Midstream. E11: the EI ``Import
share'' (a self-sufficiency ratio) and the EV ``import\_share'' (a net-export ratio) are renamed for
what they measure, and non-producers now tie exactly in Energy Imports. E12: an economy without its
own EI production row (EI counts it in a regional residual) was read as producing nothing, so New
Zealand counted as importing all its gas; production then comes from EIA International Energy
Statistics \citep{eia2026international}, flagged and counted as observed, paired with consumption of
the same year (Estonia coal stays missing because EI counts oil shale). On the inputs of 10 October
2026 this moves 605 ES Upstream fossil cells and no allied pick.

\paragraph{Trade and Economic Opportunity (T1--T13).} T1: the ES and EO trade themes were identical
tables and the ES deficit term was never computed; ES Trade is now concentration, basket share and
net basket imports over GDP, EO Trade share, RCA \citep{balassa1965trade}, export size, feasibility
and concentration (earlier--corrected Spearman 0.54--0.91 across cells). T2: export size and
feasibility were ranked within reporter, so Moldova and Samoa ranked beside China in Solar
Midstream. T3: market share was the export-weighted mean of product-level world shares; the Congo
had 0.862 against a basket share of 0.0074 in Batteries Midstream. T4: ores, slag and scrap were
filed under Batteries Midstream, untreated softwood logs under three grid and renewable baskets, and
16 keys appeared twice. T5: 73 crosswalk codes had no HS92 counterpart, so the Atlas EV Midstream
basket contained no cars; 19 predecessor codes bridge them. T6: the levelised cost entered Cost
Competitiveness four times (0.78 of the category's weight where quoted). T7: readiness was written for
a pseudo-country and reached real economies only by imputation. T8: Technology Demand listed
composites beside their own sub-indices and ranked a zero addressable market as the smallest. T9: the
export-concentration index was ranked across every basket, so unscored baskets (Semiconductors)
moved scored cells; it is now ranked over the headline cells. T10: the Comtrade energy-trade file
lacked 60 basket codes and 13 reporters, among them Qatar's LNG (USD 56.9bn in 2024); it is rebuilt
from the World rows of the all-reporter dyad pull. T11--T13: duplicate WEO and BNEF demand keys,
IEA 2035 manufacturing shares taken without the ``Other'' row (China's 2035 solar share 0.90 rather
than 0.80), zero lending-rate stubs (New Zealand 0.0 percent, now 4.76) and mislabelled IMF price
units \citep{imf2026pcps}; only the lending rates move an index value.

\paragraph{Policy and Strategic indices (P1--P9).} P1: the NIPO module's jurisdiction weights
overwrote the IEA module's, so international IEA policies were weighted 0.20 instead of 0.66. P2:
Upstream stage tags matched substrings (``ore'' in ``more'', ``Korea'', ``Singapore''), tagging 36.9
percent of IEA policies Upstream; whole-word rules tag 12.5 percent. P3: name mismatches left South
Korea, the United Arab Emirates (coded ``UAE''), France, Italy and Denmark without a Policy index;
EU members without a national CAT rating now take the EU's \citep{cat2026ratings}. P4: the policy
block had stopped before writing its table since 6~February 2026, so the September Strategic Index
combined September pillars with February policy components. P5: NIPO's default as-of date was the
latest announcement or implementation date (1~October 2028 in every 2026 extract), counting
scheduled measures as in force; it is now the extract's last announcement, capped at the release
month and run date. NIPO is no longer in the composite; the robustness variant uses the October 2026
extract as of 5~October 2026. P6: the strategic table overwrote the security column with $1-\ES$; it
now carries security and need under their own names (no value changes). P7: missing pillar or policy
values were filled with country means (51 of 571 shipped cells); a cell now needs all three. P8:
sectors with a constant pillar were not detected; Nuclear Upstream, constant at zero, was the top
pick of 8 of 18 shipped allies. P9: the staged oil lifecycle emissions (720 g CO$_2$eq/kWh) appear in
no IPCC table; the SRREN median of 840 \citep{moomaw2011annex} replaces it next to the AR5 values of
the other technologies \citep{schlomer2014annex}, so the oil weight falls from 0.12 to 0 and coal's
rises from 0 to 0.02.

\paragraph{Bilateral indices (B1--B5).} B1: grouped \texttt{transmute()} calls, name keys and a
self-crossed backbone duplicated the Friendshore table (17.5 million rows for 1.44 million cells),
so a dyad's best-five mean weighted its top cells 2, 2 and 1. B2: \texttt{mean(a:b)} returned
\texttt{a}, so outward investment was the rank of the level only; 99 of 259 IMF counterpart names did
not match, so China and the Netherlands scored zero from every reporter. B3: development potential,
a partner-level construct, was stored per reporter with donor fan-out up to 38 rows, with 56
aggregates and mis-oriented WDI series. B4: a high-income partner's undefined development score kept
its weight, capping its PSI at 0.8 and placing the United States 185th of 218. B5: an opportunity
column named for security held the partner's need, and the tables covered 18 to 20 pseudo-technologies
rather than the eleven headline ones.

\paragraph{Network (N1--N7).} These defects concerned the network design that the paper no longer
uses; Appendix~\ref{app:corrections} records them, with the design itself, in one paragraph.
